%============ Chapter2.tex — Background and Literature Review ===========
\chapter{Background and Literature Review}

\section{Overview}

This chapter provides the scientific foundation on which the project is
built. Section~2.2 summarizes the clinical background of Parkinson's
Disease (PD). Section~2.3 explains how the human voice is produced and
how PD alters it. Section~2.4 defines the acoustic measurements used to
quantify those alterations, together with the software tools that
compute them. Section~2.5 introduces the classical Machine Learning (ML)
models employed in this project. Section~2.6 examines the validation
problem that affects much of the published literature in this field.
Section~2.7 reviews related work and the public datasets used, and
Section~2.8 states the research gap that this project addresses.

\section{Parkinson's Disease}

PD is a chronic, progressive neurodegenerative disorder caused by the
loss of dopamine-producing neurons, predominantly in the substantia
nigra pars compacta, a region of the midbrain involved in the control of
movement. Post-mortem morphometric analysis indicates that this loss is
regionally selective and that a substantial proportion of the affected
neurons has already degenerated by the time motor symptoms appear
\cite{fearnley1991}. The cardinal motor manifestations are bradykinesia
(slowness of movement), muscular rigidity, rest tremor, and, in later
stages, postural instability; the disease also produces a wide range of
non-motor symptoms \cite{postuma2015}.

The burden of the disease is substantial and increasing. The Global
Burden of Disease study estimated that about 6.1 million individuals
were living with PD worldwide in 2016, compared with 2.5 million in
1990, an increase driven mainly by the growing number of older people
\cite{gbd2018}. There is at present no cure and no routine laboratory
test; the diagnosis is clinical, based on the motor examination and
supporting criteria formalized by the Movement Disorder Society (MDS)
\cite{postuma2015}.

Two aspects of this clinical picture motivate the present project.
First, because the diagnosis depends on visible motor signs, it is
typically made only after the underlying neuronal loss is well advanced
\cite{fearnley1991}. Second, the muscles of respiration, the larynx,
and the articulators are affected by the same motor impairment as the
limbs, so measurable changes in the voice can accompany --- and in some
reported cases precede --- the visible signs \cite{harel2004,rusz2011}.
An inexpensive measurement channel that reflects early motor
impairment is therefore of practical interest as a screening-support
signal, provided its output is interpreted cautiously and never treated
as a diagnosis.

\section{Speech Production and Hypokinetic Dysarthria}

\subsection{The Source--Filter Model of Voice Production}

Human voice production is conventionally described by the source--filter
model, illustrated in Figure~\ref{fig:sourcefilter}. Three stages
cooperate:

\begin{enumerate}
  \item \textbf{Power source.} The lungs supply a steady stream of
        pressurized air through the trachea.
  \item \textbf{Sound source.} The two vocal folds of the larynx are
        brought together across the airway. Subglottal pressure forces
        them apart, and their elasticity, together with aerodynamic
        forces, snaps them shut again. This cycle repeats quasi-%
        periodically, releasing a train of air puffs. The repetition
        rate of this cycle is perceived as the pitch of the voice and is
        called the Fundamental Frequency (F0); typical adult values lie
        roughly between 85 and 180~Hz for men and 165 and 255~Hz for
        women.
  \item \textbf{Filter.} The pharynx, oral cavity, tongue, teeth, and
        lips form a resonating tube whose shape changes continuously
        during speech. This vocal-tract filter emphasizes certain
        frequency regions and attenuates others, shaping the buzz of
        the source into distinguishable vowels and consonants.
\end{enumerate}

\begin{figure}[htb]
\centering
\begin{tikzpicture}[
  node distance=8mm and 17mm,
  stage/.style={draw, rounded corners=2pt, minimum width=3.0cm,
                minimum height=1.5cm, align=center, thick},
  lbl/.style={font=\small\itshape, align=center},
  arr/.style={-{Stealth[length=3mm]}, thick}
]
\node[stage] (lungs) {Lungs\\[-2pt]{\small (power source)}};
\node[stage, right=of lungs] (folds) {Vocal folds\\[-2pt]{\small (sound source)}};
\node[stage, right=of folds] (tract) {Vocal tract\\[-2pt]{\small (filter)}};
\node[right=9mm of tract, lbl] (speech) {Speech\\sound};
\draw[arr] (lungs) -- node[above=1pt, font=\footnotesize] {airflow} (folds);
\draw[arr] (folds) -- node[above=1pt, font=\footnotesize, align=center] {glottal\\pulses} (tract);
\draw[arr] (tract) -- (speech);
\node[lbl, below=4mm of lungs] {respiratory\\muscles};
\node[lbl, below=4mm of folds] {F0, jitter,\\shimmer, HNR, CPPS};
\node[lbl, below=4mm of tract] {MFCC,\\delta-MFCC};
\end{tikzpicture}
\caption{The source--filter model of voice production. The lower row
indicates which acoustic feature families used in this project
characterize each stage.}
\label{fig:sourcefilter}
\end{figure}

Two properties of this system matter for the measurements that follow.
The rate of vocal-fold vibration determines F0, and because typical F0
ranges differ between adult men and women, absolute pitch values partly
encode the speaker's sex --- a potential confounder examined explicitly
in Chapter~4. The \emph{regularity} of the vibration determines
perceived voice quality: in a healthy voice, successive cycles are
nearly identical, while irregular timing or amplitude is heard as
roughness, and incomplete closure of the folds is heard as breathiness.

\subsection{Hypokinetic Dysarthria in Parkinson's Disease}

Dysarthria denotes a group of speech disorders caused by impaired
neuromuscular control of the speech apparatus. In their classical
perceptual study, Darley, Aronson, and Brown identified distinct
dysarthria patterns corresponding to different neurological lesions and
associated parkinsonism with \emph{hypokinetic} dysarthria
\cite{darley1969}. Its perceptual hallmarks are reduced loudness,
monotone pitch (reduced F0 variability), breathy or hoarse voice
quality, imprecise consonant articulation, and disturbed speech timing.

These perceptual signs have measurable acoustic correlates, and their
prevalence is high. In a sample of 200 PD patients, speech or voice
impairment was present in approximately 90\% \cite{ho1999}.
Quantitative acoustic analysis of 46 early, untreated PD patients found
some form of vocal impairment in about 78\%, with measures of F0
variability among the most discriminative \cite{rusz2011}. Timing is
affected as well: compared with healthy speakers reading the same text,
PD patients showed altered pause behaviour and a tendency to accelerate
their speech rate \cite{skodda2008}. A longitudinal case study further
suggested that reduced F0 variability can be detected in recordings
made about five years before clinical diagnosis \cite{harel2004}.

These findings define the measurement targets for this project: pitch
variability, cycle-to-cycle regularity, the balance of periodic and
noise energy, articulatory spectral dynamics, and pause behaviour.

\section{Acoustic Measurement of the Voice}

This section defines the acoustic quantities used in the project. All
perturbation and cepstral measures are computed with the Praat speech
analysis system \cite{boersma2001} through its Python interface
Parselmouth \cite{jadoul2018}; spectral features are computed with the
librosa library \cite{mcfee2015}.

\subsection{Fundamental Frequency Statistics}

F0 is estimated over short overlapping analysis frames using the
autocorrelation method of Boersma \cite{boersma1993}, which locates the
dominant periodicity of the waveform within a configured search range
(75--500~Hz in this project, covering adult male and female voices).
Frames without detectable periodicity (silences, unvoiced consonants)
are excluded. From the resulting F0 track, summary statistics are
computed: mean, median, standard deviation, minimum, maximum, range,
and the voiced fraction (the proportion of frames in which voicing was
detected). Given the monotone character of hypokinetic dysarthria, the
dispersion statistics (standard deviation and range) are the most
clinically meaningful \cite{rusz2011,harel2004}.

\subsection{Jitter}

Jitter quantifies the cycle-to-cycle instability of vocal-fold
\emph{timing}. Let $T_i$ denote the duration of the $i$-th glottal
period and $N$ the number of periods. Local jitter is the mean absolute
difference between consecutive periods, normalized by the mean period:
\begin{equation}
\mathrm{jitter}_{\mathrm{local}} =
\frac{\dfrac{1}{N-1}\displaystyle\sum_{i=1}^{N-1}\left|T_i - T_{i+1}\right|}
     {\dfrac{1}{N}\displaystyle\sum_{i=1}^{N} T_i}.
\label{eq:jitter}
\end{equation}
Variants used in this project average over small neighbourhoods of
periods to reduce sensitivity to isolated measurement errors: absolute
local jitter (the numerator of Eq.~\ref{eq:jitter}, in seconds), the
Relative Average Perturbation (RAP, three-period neighbourhood), and
the five-period Perturbation Quotient (PPQ5). In healthy sustained
phonation, local jitter is typically below about 1\%; impaired
neuromuscular control raises it.

\subsection{Shimmer}

Shimmer applies the same idea to cycle \emph{amplitudes} $A_i$:
\begin{equation}
\mathrm{shimmer}_{\mathrm{local}} =
\frac{\dfrac{1}{N-1}\displaystyle\sum_{i=1}^{N-1}\left|A_i - A_{i+1}\right|}
     {\dfrac{1}{N}\displaystyle\sum_{i=1}^{N} A_i}.
\label{eq:shimmer}
\end{equation}
The project uses local shimmer, its decibel form, and the three-, five-,
and eleven-period Amplitude Perturbation Quotients (APQ3, APQ5, APQ11).
Weak or inconsistent vocal-fold closure produces uneven pulse strength
and therefore raised shimmer. Because shimmer is normalized by the mean
amplitude, a uniform change of recording gain does not affect it.

\subsection{Harmonics-to-Noise Ratio}

The voice signal is a mixture of a periodic component (the harmonics of
F0) and an aperiodic noise component caused by turbulent airflow. The
Harmonics-to-Noise Ratio (HNR) expresses their energy ratio in
decibels. Using the normalized autocorrelation approach of Boersma, if
$r_{\max}$ is the height of the normalized autocorrelation peak at the
detected period (the fraction of signal power that is periodic), then
\begin{equation}
\mathrm{HNR} = 10\log_{10}\frac{r_{\max}}{1-r_{\max}}\;\text{dB}.
\label{eq:hnr}
\end{equation}
A clear voice may exceed 20~dB, while breathy voices score much lower
\cite{boersma1993}. Incomplete glottal closure --- a recognized feature
of hypokinetic dysarthria --- lowers HNR.

\subsection{Smoothed Cepstral Peak Prominence}

Jitter, shimmer, and HNR presuppose that individual glottal cycles can
be identified reliably, which holds for sustained vowels but only
approximately for continuous speech. The Smoothed Cepstral Peak
Prominence (CPPS) avoids this requirement. The cepstrum is obtained by
taking a second spectral transform of the logarithmic power spectrum;
a strongly periodic voice, whose spectrum contains regularly spaced
harmonics, produces a sharp cepstral peak at the quefrency of the
glottal period. CPPS measures how far this peak rises above the
regression line fitted to the smoothed cepstrum, in decibels. A
meta-analysis of acoustic correlates of perceived voice quality found
cepstral peak prominence measures to be among the strongest and most
robust correlates of overall dysphonia, in both sustained vowels and
continuous speech \cite{maryn2009}. Breathier, noisier voices produce
weaker cepstral peaks and therefore lower CPPS.

\subsection{Pause and Timing Statistics}

Speech timing is quantified here with an energy-based segmentation of
the recording into speech stretches and silences. Gaps of at least
0.2~s are counted as pauses, from which four features are derived:
pauses per minute, mean pause duration, the fraction of total time
spent in pauses, and the mean duration of continuous speech stretches.
The clinical motivation follows from the altered pause behaviour
documented in PD reading tasks \cite{skodda2008}. These features have
the additional practical advantage of depending on timing rather than
on spectral content, which makes them comparatively robust to changes
of microphone and recording channel.

\subsection{Mel-Frequency Cepstral Coefficients and Their Dynamics}

The Mel-Frequency Cepstral Coefficients (MFCC) provide a compact
description of the short-term spectral envelope --- effectively, of the
vocal-tract filter configuration. The signal is divided into short
frames; for each frame the power spectrum is computed, integrated into
overlapping bands spaced on the mel scale,
\begin{equation}
m(f) = 2595\,\log_{10}\!\left(1+\frac{f}{700}\right),
\label{eq:mel}
\end{equation}
which approximates the frequency resolution of human hearing; the band
energies are logarithmized; and a discrete cosine transform yields the
coefficients, of which the first 13 are retained here, following the
representation introduced by Davis and Mermelstein \cite{davis1980}.
Delta-MFCC coefficients are the frame-to-frame time derivatives of the
MFCCs and describe how quickly the spectral envelope changes ---
a proxy for articulatory movement speed. In this project each
coefficient trajectory is summarized by its mean and standard deviation
over the analysed audio, yielding 26 MFCC and 26 delta-MFCC features.
Because MFCCs describe the spectrum in detail, they are also sensitive
to the microphone, room, and language --- a property with consequences
examined in the external validation (Chapter~4).

\section{Classical Machine Learning Models}

Four classifier families, all implemented in the scikit-learn library
\cite{pedregosa2011}, are used in this project, together with a
majority-class baseline that always predicts the more frequent class
and therefore represents zero learned knowledge.

\textbf{Logistic regression} models the probability of class
membership as a logistic function of a weighted sum of the features.
It produces a linear decision boundary, and the learned weights
indicate each feature's direction and strength of influence, making it
the most directly interpretable model considered.

\textbf{The Support Vector Machine (SVM)} \cite{cortes1995} seeks the
decision boundary that separates the classes with the widest possible
margin, which favours good generalization from limited data. With the
Radial Basis Function (RBF) kernel, the data are implicitly mapped into
a higher-dimensional space, allowing the boundary to be curved while
the margin principle is retained.

\textbf{The Random Forest (RF)} \cite{breiman2001} is an ensemble of
decision trees, each trained on a bootstrap sample of the training data
with random feature subsets considered at each split; the forest
predicts by aggregating the votes of its trees. RFs handle features of
mixed scales and nonlinear interactions naturally, resist overfitting
through averaging, and provide a ranking of feature importance ---
a property used in Chapter~4 to interpret the final model.

\textbf{The Multilayer Perceptron (MLP)} is a small feed-forward neural
network trained by gradient descent. It is included for completeness as
the most flexible of the four families, but flexibility carries risk
with only 37 subjects, and its behaviour under a confounder ablation
(Chapter~4) illustrates exactly this danger.

The rationale for restricting the project to such classical models
rather than deep learning is threefold: the sample size (37 subjects)
is far below what deep networks require; the biomedical setting demands
interpretable evidence of \emph{what} the model uses; and the project's
scientific focus is on honest validation, which is clearer to
demonstrate with simple, well-understood models.

\section{Validation Methodology and the Data-Leakage Problem}

The value of any reported accuracy depends entirely on how the model
was tested. When a dataset contains several recordings --- or several
excerpts of one recording --- from the same person, a random split of
those items between training and test sets allows the model to
encounter each test speaker during training. The model can then succeed
by \emph{recognizing the speaker or the recording}, rather than the
disease, and the measured performance no longer predicts behaviour on
new people. This phenomenon is a form of data leakage, and in the
clinical machine learning literature the distinction is drawn between
\emph{record-wise} and \emph{subject-wise} cross-validation: Saeb et
al.\ demonstrated that record-wise validation can dramatically
overestimate diagnostic accuracy and argued that the validation scheme
must approximate the intended use case --- diagnosing \emph{new}
individuals \cite{saeb2017}.

The correct procedure is grouped, subject-independent validation: all
data from one person is assigned as a block to either the training or
the test side of every split. Where class proportions should also be
preserved across folds, stratified grouped K-fold splitting is used.
This project enforces subject independence mechanically --- a runtime
assertion halts execution if any subject ever appears on both sides of
a split --- and additionally demonstrates the size of the leakage
effect empirically on its own data, by evaluating the same pipeline
under a deliberately wrong chunk-level split (Chapter~4).

A related, subtler leakage channel concerns preprocessing statistics:
if imputation values or feature scaling parameters are estimated on the
full dataset before splitting, information about the test data enters
training. The methodology of Chapter~3 therefore fits all such
operations inside each training fold only.

\section{Related Work and Datasets}

Research on voice-based PD assessment spans two broad recording
paradigms. The first uses \emph{sustained vowels}, in which the speaker
phonates a steady ``ah''; this isolates phonatory function and suits
the classical perturbation measures. Little et al.\ combined dysphonia
measures of sustained vowels with an SVM classifier to separate PD
speakers from controls \cite{little2009}, and Tsanas et al.\ extended
this line with a set of 132 dysphonia measures and feature selection,
reporting almost 99\% classification accuracy on sustained-vowel data
\cite{tsanas2012}. Results of this magnitude, here as elsewhere in the
field, must be read together with the validation-design question
discussed in Section~2.6: when several samples of the same speaker can
fall on both sides of a random split, part of the measured performance
may reflect speaker recognition rather than disease detection
\cite{saeb2017}. The second paradigm uses \emph{connected
speech} (read text or spontaneous conversation), which engages
articulation and prosody in addition to phonation and is closer to
natural speaking conditions \cite{rusz2011,skodda2008}; it is the
paradigm of the present project.

Two public connected-speech datasets are used in this work. The
MDVR-KCL corpus \cite{jaeger2019} contains smartphone recordings of
read text and spontaneous dialogue from 37 English-speaking
participants (21 healthy controls and 16 PD patients) and serves as the
training corpus; its detailed inspection appears in Chapter~4. The
Italian Parkinson's Voice and Speech database, collected by Dimauro et
al.\ in the context of speech-intelligibility assessment, contains
recordings of 65 Italian speakers in young healthy, elderly healthy,
and PD groups \cite{dimauro2017}. The publicly mirrored copy of this
database used in the present work comprises 61 of those speakers
(15 young healthy controls, 22 elderly healthy controls, and 24 PD
patients; Chapter~4), and serves here exclusively as an external,
never-trained-on test set.

Published classification studies on MDVR-KCL and similar corpora
frequently report accuracies above 95\%. In light of the record-wise
leakage mechanism described above \cite{saeb2017}, and of the
demonstration in Chapter~4 that a chunk-level random split inflates
this project's own balanced accuracy from 0.775 to 0.909 without any
change to the model, such figures should be interpreted with caution
unless subject-independent validation is stated explicitly.

\section{Research Gap and Project Contribution}

From the literature reviewed above, three gaps emerge. First, many
reported results rest on validation schemes that do not separate
speakers, leaving their practical meaning uncertain \cite{saeb2017}. Second, cross-dataset evaluation --- testing a trained
system on an entirely different corpus, language, and recording channel
--- is rare, although it is the most direct measure of the
generalization that a screening application would require. Third,
reported systems seldom address how their output should be presented
given the uncertainty involved; a bare class label overstates what a
small-corpus model knows.

This project addresses all three gaps within an undergraduate scope:
it enforces subject-independent validation mechanically and quantifies
the leakage effect on its own data; it performs an unadapted external
validation on a second corpus in a different language; and it delivers
a prototype whose output includes an explicit inconclusive band,
recording-condition warnings, and mandatory non-diagnostic wording.
The complete methodology follows in Chapter~3.
